%--------------------------
\begin{frame}{Question 6}

\begin{block}{}
La question de l’unité de l’invention :

\begin{itemize}
  \item[$\bullet$] a. n’est pas examinée par l’office récepteur lors de la vérification des taxes acquittées. 
  \item[$\bullet$] b. n’est pas examinée par l’administration chargée de la recherche internationale après 
  l’établissement du rapport de recherche internationale. 
  \item[$\bullet$] c. peut être examinée par l’administration chargée de l’examen préliminaire 
  international même si l’administration chargée de la recherche internationale ne l’a 
  pas examinée. 
  \item[$\circ$] d. Aucune des réponses a, b ou c n’est correcte. 
\end{itemize}

\end{block}

\end{frame}

%--------------------------
\begin{frame}{Question 6}

\begin{itemize}
  \item[$\bullet$] \textbf{a.} \; Correct : l’office récepteur vérifie les taxes formelles
  mais ne procède pas à un examen de l’unité de l’invention (\pctrule{13}).

  \item[$\bullet$] \textbf{b.} \; Correct : l’examen de l’unité par l’ISA intervient pendant la recherche,
  pas après l’établissement du rapport de recherche internationale
  (\pctart{17}.3.a)).

  \item[$\bullet$] \textbf{c.} \; Correct : l’IPEA peut examiner l’unité indépendamment
  de l’ISA et inviter au paiement de taxes additionnelles
  (\pctart{34}.3.a)).

  \item[$\circ$] \textbf{d.} \; Faux puisque \textbf{a}, \textbf{b} et \textbf{c} sont correctes.
\end{itemize}

\end{frame}
